\section{Introduction}

How much can two-sided midpoint growth differ from one-sided asymptotic
convexity?  We answer this quantitatively for certain real subspaces of
$L^1$.  The norm is fixed throughout the computation: it is the inherited
$L^1$ norm.  The obstruction to AUC renorming will require a separate
argument that applies to every equivalent norm.

For an infinite-dimensional real Banach space $X$, write $B_X$ and $S_X$
for its closed unit ball and unit sphere, and
$\operatorname{cof}(X)$ for its closed linear subspaces of finite
codimension.  Given a norm $N$, a vector $x$ with $N(x)=1$, and $t>0$, set
\begin{align*}
 H_N(x,t)&=\sup_{F\in\operatorname{cof}(X)}
   \inf_{\substack{y\in F\\N(y)\ge1}}
   \left(\frac{N(x+ty)+N(x-ty)}2-1\right),\\
 D_N(x,t)&=\sup_{F\in\operatorname{cof}(X)}
   \inf_{\substack{y\in F\\N(y)=1}}\bigl(N(x+ty)-1\bigr).
\end{align*}
The averaged midpoint modulus $\widehat\delta_N(t)$ and the one-sided
modulus $\overline\delta_N(t)$ are their respective infima over unit
centers.  Thus a subspace of finite codimension is chosen for each center
before the perturbing direction is chosen.  The norm $N$ is
\emph{asymptotically uniformly convex} (AUC) if
$\overline\delta_N(t)>0$ for every $t>0$.

For each fixed $y$, both the average and the maximum of $N(x+ry)$ and
$N(x-ry)$ are even convex functions of $r$.  Thus unit directions give
the same inner infimum in $H_N$ and in its maximum-based counterpart.  Let $m_N$ be the modulus obtained by replacing the average in $H_N$ by
the maximum and then taking the infimum over unit centers.  It satisfies
$\tfrac12m_N(t)\le\widehat\delta_N(t)\le m_N(t)$.
Indeed, the upper bound is immediate; for the lower bound, intersect a
candidate subspace with the kernel of a norm-one functional norming $x$.
Both perturbed norms are then at least one, and the average of their
excesses is at least half their maximum excess.  Taking the infima and
suprema gives the comparison.  Positivity of the averaged modulus is
therefore the usual asymptotic midpoint uniform convexity (AMUC).

Dilworth, Kutzarova, Randrianarivony, Revalski, and Zhivkov introduced
AMUC and asked whether AMUC and AUC are equivalent up to renorming
\cite[Definition~2.2 and the question in Section~5]{DKRRZ2016}.
Baudier answered this question negatively by proving
AMUC for subspaces of $L^1$ whose unit balls are relatively compact in
measure and applying this to the Daugavet example of Kadets and Werner
\cite[Theorem~1 and Corollary~1]{Baudier2026}.  We calculate both moduli
exactly under these two geometric hypotheses.  Baudier's proof uses
the $L^1$ splitting identity of Br\'ezis and Lieb
\cite[Theorem~1]{BrezisLieb1983} along almost-everywhere convergent
subsequences.  Our lower estimate instead uses a finite net of pointwise
truncations and the sign functionals attached to that net.

Recall that a Banach space has the \emph{Daugavet property} if
$\|I+T\|=1+\|T\|$ for every rank-one bounded operator $T$ on the space.
For functions on a probability space, convergence in measure is induced
by the metric
\begin{equation}\label{eq:measure-metric}
 d(f,g)=\inf\{a>0:\mathbb P(|f-g|>a)<a\}.
\end{equation}
Total boundedness for this metric means that every positive radius
admits a finite covering by balls of that radius. The following theorem
combines this compactness assumption with the Daugavet property.

\begin{theorem}\label{thm:exact}
Let $E$ be an infinite-dimensional closed real subspace of
$L^1(\Omega,\mathbb P)$, where $\mathbb P$ is a probability measure.
Suppose its closed unit ball is totally bounded for convergence in
measure and its inherited norm has the Daugavet property.  Then every
$x\in S_E$ and $t>0$ satisfy
\[
 H_E(x,t)=\max\{t/2,t-1\},\qquad
 D_E(x,t)=\max\{0,t-2\}.
\]
Consequently these are also the formulas for
$\widehat\delta_E(t)$ and $\overline\delta_E(t)$.  The space $E$ admits
no equivalent AUC norm.
\end{theorem}

In particular, the averaged modulus has linear growth near the origin,
whereas the one-sided modulus remains zero up to $t=2$.
Baudier also gives a universal linear lower bound for the maximum-based
modulus of the Kadets--Werner space
\cite[final paragraph]{Baudier2026}.  Our refinement determines the exact
averaged and one-sided curves at every unit center under the two
geometric hypotheses.
Figure~\ref{fig:moduli} displays the two curves.
\begin{figure}[H]
\centering
\begin{tikzpicture}[x=1.65cm,y=1.1cm,>=Stealth,font=\small]
 \draw[->] (0,0)--(4.65,0) node[right] {$t$};
 \draw[->] (0,0)--(0,3.65) node[above] {modulus};
 \foreach \x in {1,2,3,4} {
   \draw (\x,.04)--(\x,-.04) node[below] {$\x$};}
 \foreach \y in {1,2,3} {
   \draw (.04,\y)--(-.04,\y) node[left] {$\y$};}
 \draw[gray!50,dashed] (2,0)--(2,1);
 \draw[blue!70!black,very thick] (0,0)--(2,1)--(4.3,3.3);
 \draw[orange!80!black,thick,dashed] (0,0)--(2,0)--(4.3,2.3);
 \node[blue!70!black,anchor=west] at (.2,3.25)
   {$\whdelta_E(t)=\max\{t/2,t-1\}$};
 \node[orange!80!black,anchor=west] at (.2,2.7)
   {$\bardelta_E(t)=\max\{0,t-2\}$};
\end{tikzpicture}
\caption{The exact moduli of the original norm on $E$.
The averaged midpoint modulus is positive for every positive radius,
whereas the one-sided modulus vanishes through radius $2$.
Theorem~\ref{thm:exact} gives both equalities at each unit center.
These curves do not describe other equivalent norms.}
\label{fig:moduli}
\end{figure}


The existence of such spaces has an older history.  Bourgain and
Rosenthal constructed an $L^1$ subspace whose unit ball is relatively
compact in measure although the space fails the Radon--Nikodym property
\cite{BourgainRosenthal1980}.  Kadets and Werner refined its
finite-dimensional enlargement to obtain the Daugavet property
\cite[Lemma~2.4 and Theorem~2.5]{KW2004}. We give a complete realization of
their enlargement and limiting scheme in Section~\ref{kw:section}, so
the existence of an example and the calculation of its moduli form one
self-contained argument.

The two assumptions enter different parts of the proof.
In Section~\ref{sec:clips}, pointwise clipping at height $|x|/t$ turns the
measure-precompact ball into a norm-compact set.  Finitely many sign
functionals then force at least half the $L^1$ mass of every unit
direction in their common kernel to remain outside the clip, up to an
arbitrarily small error. This gives the lower
averaged bound.  Section~\ref{sec:weak-neighborhoods} develops the other
input: every weak neighborhood of a unit vector contains points almost
at distance two from it.  This is the needed case of Shvydkoy's
weak-open characterization of the Daugavet property
\cite[Lemma~3]{Shvydkoy2000}, obtained from the slice characterization
of Kadets, Shvidkoy, Sirotkin, and Werner \cite[Lemma~2.2]{KSSW2000}.
In Section~\ref{sec:curves}, correcting a finite-dimensional quotient
coordinate turns those almost diametral differences into directions in
any prescribed subspace of finite codimension, giving both upper curves.
Section~\ref{sec:renorming} uses the same weak-neighborhood
property at a near-maximizer of an arbitrary equivalent norm to rule
out AUC renorming. Section~\ref{kw:section} then constructs an increasing
sequence of finite-dimensional spaces whose closed union has both
required properties. At each step, multiplying selected old unit
vectors by independent positive variables creates new vectors almost
at distance two from every old unit vector, while averages of the new
vectors recover the selected ones. A uniform estimate places the
entire enlarged unit ball close in measure to the old ball. Summable
errors preserve total boundedness through the limit, and the averaging
property gives the Daugavet slice condition.
